\phantomsection % Extrem wichtig! Setzt den "Anker" für das Inhaltsverzeichnis exakt auf diese Seite.
\section*{Vorwort}
\addcontentsline{toc}{section}{Vorwort}

%\addcontentsline{toc}{section}{Hinweise zur Arbeit mit diesem Skript} % Falls es ins Inhaltsverzeichnis soll

Dieses Dokument ist als aktiver Lesetext konzipiert. Es richtet sich an Lehrpersonen, die nach spannenden, tiefergehenden Anschauungsbeispielen für den Geometrieunterricht suchen, ebenso wie an interessierte Schülerinnen und Schüler oder Studierende, die sich das Thema im Selbststudium erarbeiten möchten.

Um die teils anspruchsvolle Theorie, insbesondere die formale Matrixalgebra und die gruppentheoretischen Beweise, greifbarer zu machen, sind in den Text regelmäßig Aufgaben eingestreut. Diese Übungen sind kein bloßes Anhängsel, sondern ein essenzieller Bestandteil des didaktischen Konzepts. Sie laden dazu ein, die abstrakten mathematischen Konzepte direkt an konkreten, oft stark visuellen Beispielen auszuprobieren und selbst zu verifizieren. 

Lernenden wird empfohlen, diese Aufgaben direkt während der Lektüre zu bearbeiten, um das eigene Verständnis Schritt für Schritt zu überprüfen. Die Theorie baut oftmals auf den Erkenntnissen dieser Übungen auf. Zur Kontrolle finden sich die vollständigen, ausführlichen Lösungswege am Ende des Dokuments. 
